% !TeX root = ../final-report.tex
\section{Introduction}
Replicating video segments across storage nodes allows a viewer to switch
sources when one node fails. Continuing playback, however, does not restore
the missing replica. Until a replacement is installed, another holder failure
can make a segment unavailable to new requests. Repair closes this window by
copying data from a surviving holder, but that holder must also serve viewers.
With limited upload capacity, protecting future availability can therefore
interfere with present playback.

We study this tension in a cooperative video storage prototype with dedicated
storage Peers, independent viewers, and a Raft-replicated Controller service.
Peers contribute disk space and upload capacity; viewers do not upload.
Controllers commit manifests, replica registrations, and repair tasks, while
video bytes remain outside the consensus log. The system requires two
confirmed copies of every segment before publication and repairs published
segments when fewer than two verified, reachable copies remain.

Our central question is: \emph{how should repair timing and bandwidth limits
balance playback continuity against replica recovery on heterogeneous,
upload-constrained Peers?} We examine whether repair adds interference beyond
the initial failure, how a grace interval trades avoided copying against
reduced redundancy, and whether a source-local load-aware allowance improves
on fixed repair caps. We also test the consequences of a second holder
failure and recovery of application state after Controller failures.

The project makes three scoped contributions. First, it implements a working
storage and playback system with durable, verifiable objects, quorum-backed
publication, resumable repair tasks, and browser playback. Second, it provides
shared, fixed-cap, and load-aware repair controls under an unchanged total
source upload budget, together with paired fault experiments and classified
traffic accounting. Third, it characterizes their trade-offs: repair can add
stalls under contention, longer grace can avoid short-outage copies while
extending exposure, and load-aware control shows a workload-specific
playback/recovery compromise rather than a consistent advantage.

The contribution is an implemented prototype and an exploratory empirical
study. Performance experiments transfer real HTTP payloads but simulate
media consumption; browser decoding is validated separately. The final scope
uses replication, without the erasure coding, popularity learning, or tiered
storage envisioned in the earlier proposal.
